\subsection{Threats to validity and limitations}
\label{sec:limitations}

\subsubsection{Construct and data validity}

NSRDB GHI is a satellite-informed physical-model estimate
\cite{sengupta2018}, not an on-site pyranometer record. Factory coordinates
select product grid cells, so neither the target nor its error quantifies the
irradiance at a particular roof surface. The purposive set of ten
factory-associated locations is also not a sample of factory electricity
demand, photovoltaic output, roof area, or operational constraints. The
industrial framing therefore identifies geographically relevant forecasting
sites; it does not establish plant-level energy savings.

The matched input contains only GHI available at the forecast origin and
site identity. This restriction controls predictor availability among learned
models, but excludes potentially useful weather information, such as cloud
motion, numerical weather predictions, and forecasts of aerosols and
temperature. A site embedding can also encode persistent
location-specific behavior and does not demonstrate transfer to an unseen
factory. Fixed local offsets produce regular hourly grids but do not reproduce
daylight-saving clock changes. Moreover, the daily/monthly archive was
retrieved independently of the hourly archive. The comparison across resolutions
therefore does not use controlled reaggregation of a single set of native
records.

\subsubsection{Temporal and statistical validity}

All test evidence is retrospective and restricted to 2024. It does not measure
prospective robustness under later satellite-product revisions, climate
nonstationarity, or new sites. The retrospective information set assumes
that each completed historical aggregate is immediately available; product
publication delays and operational data latency are not simulated. The
primary monthly task has only 48 pre-test refit origins per site, and the
exploratory six-month task has seven test origins per site. These small
samples limit the evidence available for fitting and evaluating monthly
models relative to the hourly and daily tasks.

Targets from adjacent multi-step origins overlap, creating strongly dependent
errors. Site-first aggregation avoids allowing that overlap to change spatial
weights, but it does not make the origins independent. The paired summaries
therefore use ten locations as descriptive units, without applying an
i.i.d. significance test to the overlapping target points. Nearby sites may
also share weather and model errors, so their summaries cannot be assumed
to be independent statistical replicates.

\subsubsection{Model-comparison validity}

TimesNet and DilatedRNN share inputs, origins, seeds, optimization limits, and
post-processing in the daily and monthly tasks, but their parameterizations and
inductive biases are not identical. The frozen configuration is a controlled
comparison rather than an exhaustive architecture-specific search. Moreover,
the hourly extension evaluates DilatedRNN using seed 42 against previously
generated comparator forecasts, whereas daily and monthly estimates use five
seeds. Comparisons across resolutions therefore reflect these protocol
differences, and seed variability cannot be assessed for the hourly task.
The reported TimesNet is a compact variant with sample-specific Fourier-period
selection; the findings do not establish a ranking of every implementation
bearing either architecture name.

Neural epoch selection minimizes normalized validation MSE, whereas the main
test ranking minimizes macro-MAE in W/m$^2$. These objectives weight errors
differently, and equal selection rules do not establish each model's best
achievable MAE. Calendar baselines also use explicit target calendar labels
and historical lookups unavailable as features to the learned models. Their
performance is a useful operational reference, but does not isolate
architecture under identical feature representations.

The neural parameter budgets also differ by task, as reported in
Table~\ref{tab:model_capacity}. Equal optimization rules do not imply equal
representational capacity, and the comparison cannot isolate whether a result
comes from recurrence, periodic convolution, parameter count, or their
interaction. Training and inference runtimes were not measured consistently
across runs, so the results do not support a hardware-efficiency ranking.

The zero floor enforces nonnegative predictions at all resolutions. The
additional hourly nighttime mask uses deterministic solar geometry, but no
common upper physical envelope is imposed. Consequently, implausibly high
daytime extrapolations remain reflected in the reported errors.

Hourly metrics retain all target hours after the common nighttime mask. Exact
nighttime zeros are part of the operational 72-hour trajectory, but they also
dilute the contribution of daylight errors to the all-hour average. The
reported hourly values must therefore be interpreted as full-trajectory
metrics. In particular, exact leads 24, 48, and 72 sample only the nighttime
23:00 intervals and cannot distinguish model skill in this experiment.
Daylight-only evaluation and complete lead profiles are needed to assess
the irradiance-producing hours separately.

\subsubsection{Contextual-data validity}

The K\"oppen--Geiger labels are climate classes extracted from cells of the
present nominal 1-km map \cite{beck2018}; they do not resolve microclimates within a cell or
conditions on a specific roof. Class boundaries are uncertain, as made explicit
by the 50\% confidence of the BMW San Luis Potos\'i assignment. Climate labels
are descriptive strata, not model inputs and not explanations of individual
forecast errors.

NASA POWER is likewise a gridded, model-based external product rather than
local rain-gauge or ceilometer evidence \cite{nasapowerdaily}. The 1~mm/day
precipitation and 80\% cloud thresholds provide reproducible post-hoc tags, but
thresholding cannot establish that weather caused a particular error.
These weather observations provide context for the selected cases, while
model rankings are based on the full aggregate evaluation.
